We check whether the maximum displacement varies linearly with the initial input. The sensor on the spigot is 1mm away from where the force is being applied  
\begin{figure}[h] \centering
    \begin{tikzpicture}[scale=0.95]
        \begin{loglogaxis}[
            % width = 8cm, height = 6cm,
            xlabel = {Input amplitude [m]}, ylabel = {Maximum displacement [m]},
            legend pos = north west,
        ]
            \addplot+[no marks] table [col sep=comma, x index = 0, y index=1]
            {Modelling/Forces/Amplitudes/base_100hz.csv};
            \addlegendentry{Spigot}
            \addplot+[ no marks] table [col sep=comma, x index = 0, y index=1,]
            {Modelling/Forces/Amplitudes/stem_100hz.csv};
            \addlegendentry{Stem}
            \addplot+[no marks] table [col sep=comma, x index = 0, y index=1,]
            {Modelling/Forces/Amplitudes/collar_100hz.csv};
            \addlegendentry{Collar}
        \end{loglogaxis}
    \end{tikzpicture}
\caption{Maximum displacement as a function of input amplitude at \SI{100}{Hz}.}
\label{fig:}
\end{figure}

\section{Displacements over time}
The simulations are made to run for 40/frequency, which was expected to be enough time to observe the input propagating through the device.

Figure \ref{fig:displacement-time-1v900} shows displacement of the same 1 hz compared to 900hz. Point is to highlight that at 1hz it's time-invariant, and 900hz begins with a phase-shift and over time they become in phase.
We attempt to quantify this shifting phase with a spectrogram and in later sections discuss simulations that run for 200/frequency.

\begin{figure}[h]\centering
    \begin{subfigure}{0.49\textwidth}
        \begin{tikzpicture}[scale=0.95]
            \begin{axis}[
                xlabel = {Time [s]}, ylabel = {Displacement [m]},
            ]
                \addplot+[ no marks] table [col sep=comma, y expr = -\thisrowno{3}]
                    {Modelling/Frequencies/csv/Job-00_1-Hz_Node BASE-1.8.csv};
                \addlegendentry{Spigot}
                \addplot+[ no marks] table [col sep=comma, y expr = \thisrowno{3}]
                    {Modelling/Frequencies/csv/Job-00_1-Hz_Node STEM-1.17.csv};
                \addlegendentry{Stem}
                \addplot+[ no marks] table [col sep=comma, y expr = -\thisrowno{3}]
                {Modelling/Frequencies/csv/Job-00_1-Hz_Node COLLAR-1.7.csv};
                \addlegendentry{Collar}
            \end{axis}
        \end{tikzpicture}
        \caption{1 Hz}
    \end{subfigure}
    \begin{subfigure}{0.49\textwidth}
        \begin{tikzpicture}[scale=0.95]
            \begin{axis}[
                xlabel = {Time [ms]}, ylabel = {Displacement [nm]},
            ]
                \addplot+[ no marks] table [col sep=comma, y expr = 1e9*\thisrowno{3}, x expr = 1e3*\thisrowno{0}]
                    {Modelling/Frequencies/csv/Job-22_900-Hz_Node BASE-1.8.csv};
                \addlegendentry{Spigot}
                \addplot+[ no marks] table [col sep=comma, y expr = 1e9*\thisrowno{3}, x expr = 1e3*\thisrowno{0}]
                    {Modelling/Frequencies/csv/Job-22_900-Hz_Node STEM-1.17.csv};
                \addlegendentry{Stem}
                \addplot+[ no marks] table [col sep=comma, y expr = 1e9*\thisrowno{3}, x expr = 1e3*\thisrowno{0}]
                {Modelling/Frequencies/csv/Job-22_900-Hz_Node COLLAR-1.7.csv};
                \addlegendentry{Collar}
            \end{axis}
        \end{tikzpicture}
        \caption{900 Hz}
    \end{subfigure}

    \caption{Displacement in the z-axis. in (a) we take the opposite value for the predominantly negative curves, i.e. spigot and collar for easier comparison. }
    \label{fig:displacement-time-1v900}
\end{figure}

Taking this same 900 Hz signal, we can look at the FFT of the signal:
\begin{figure}[h]\centering
    \begin{tikzpicture}
        \begin{axis}[
            xlabel = {Frequency [Hz]}, ylabel = {Amplitude },
            width = 14cm, height=8cm,
        ]
            \addplot+[ no marks] table [col sep=comma]
                {Modelling/Frequencies/fft_spigot_900Hz.csv};
            \addlegendentry{Spigot}
            \addplot+[ no marks] table [col sep=comma]
                {Modelling/Frequencies/fft_stem_900Hz.csv};
            \addlegendentry{Stem}
            \addplot+[ no marks] table [col sep=comma]
            {Modelling/Frequencies/fft_collar_900Hz.csv};
            \addlegendentry{Collar}
        \end{axis}
    \end{tikzpicture}
    \caption{FFT given sinusoidal input at 900hz}
    \label{fig:}
\end{figure}
We quantify the change in phase over time: